\section{MIRE: Model-Irrelevance Containment Engine}
\label{sec:mire}

\subsection{The Impossibility Result and Paradigm Shift}

The design of \mire begins with a negative result. Goldwasser and
Kim~\cite{goldwasser2022} prove that no polynomial-time algorithm can
distinguish a backdoored model from a clean model using only clean-data
evaluation. Formally, for any efficient distinguisher $D$ and any backdoor
insertion procedure $B$ that preserves clean-data accuracy:
\[
  \left|\Pr[D(M_{\text{clean}}) = 1] - \Pr[D(B(M_{\text{clean}})) = 1]\right|
  \leq \text{negl}(\lambda)
\]
where $\lambda$ is the security parameter. This is not an engineering
limitation amenable to incremental improvement---it is a \emph{mathematical
proof}. No amount of monitoring sophistication, dataset curation, or
computational budget overcomes an impossibility result. Detection has a
fundamental ceiling.

\mire responds to this impossibility by reframing the question entirely.
The conventional question---\emph{how do we detect whether a model is
backdoored?}---is the wrong question. The correct question is:
\textbf{how do we make the backdoor irrelevant?} If a backdoor activates
but accomplishes nothing---if its output is blocked, its behavioral
signature is flagged, and its architectural reach is contained---then the
distinction between a backdoored and clean model ceases to matter from a
security standpoint.

This is a paradigm shift from \emph{detection} to \emph{containment}.
Rather than attempting to identify whether a model has been compromised
(provably impossible in the general case), \mire ensures that compromise
\emph{does not translate into impact}. The six mechanisms described below
implement this containment strategy across complementary attack surfaces.

\subsection{Containment Mechanisms}

\mire comprises six containment mechanisms (M1--M6), each targeting a
distinct layer of the attack surface. No single mechanism is sufficient;
the composition of all six provides the containment guarantee.

\paragraph{M1: Output Envelope Validator.}
The first and most critical mechanism is a deterministic output validator
implemented in Rust---crucially, \emph{not} a machine learning model. The
validator enforces four classes of constraint on every model output before
it reaches any downstream consumer:

\begin{enumerate}[leftmargin=*]
\item \textbf{Schema conformance.} The output must match the expected
      structural schema (JSON fields, type constraints, length bounds).
\item \textbf{Capability bounds.} The output must not reference tools,
      resources, or actions outside the model's current capability envelope.
\item \textbf{Policy compliance.} The output must satisfy all active policy
      rules (content restrictions, rate limits, sensitivity labels).
\item \textbf{Content safety.} The output must pass deterministic content
      filters for known-harmful patterns.
\end{enumerate}

\noindent
The key insight is that a backdoored model can generate \emph{arbitrary}
output, but the validator operates on the output side of the trust boundary.
Because the validator is a simple, deterministic program with no learned
parameters, it is amenable to formal verification: one can prove that no
output violating the envelope specification passes the validator, regardless
of what the model produces. The validator's correctness does not depend on
the model's integrity.

\paragraph{M2: Canary Probes.}
Canary probes provide continuous behavioral verification without imposing
latency on the user-facing path. A battery of 50 known-answer probes is
maintained, covering four categories: factual accuracy (verifiable
ground-truth questions), safety compliance (prompts that must be refused),
instruction following (prompts with deterministic expected outputs), and
refusal behavior (edge cases near policy boundaries). Probes are injected
asynchronously alongside production traffic and evaluated off-path, adding
0\,ms to user-visible latency.

The probe set rotates every 1{,}000 requests to prevent a compromised model
from learning to recognize and selectively answer probes correctly. Probe
responses are evaluated using CUSUM (Cumulative Sum) change-point
detection~\cite{page1954}, which accumulates deviations from the expected
response distribution. An alert fires when the CUSUM statistic exceeds a
$3\sigma$ threshold, indicating a statistically significant shift in model
behavior. Because probes are drawn from a distribution unknown to the model,
a backdoor that alters model behavior on adversarial triggers will
inevitably degrade canary performance---even if the backdoor is designed to
preserve clean-data accuracy on standard benchmarks.

\paragraph{M3: Spectral Watchdog.}
The spectral watchdog applies Random Matrix Theory (RMT) to model security
monitoring---to our knowledge, the first application of RMT in this context.
The mechanism periodically samples weight matrices from the deployed model
and analyzes their eigenvalue spectra to detect unauthorized modifications.

The procedure operates in four stages. First, a random subset of weight
matrices $\{W_1, \ldots, W_k\}$ is sampled from the model. Second, singular
value decomposition (SVD) is applied to obtain the empirical spectral
distribution $\rho_{\text{emp}}(\lambda)$ for each sampled matrix. Third,
the empirical distribution is fitted against the Marchenko-Pastur law, which
characterizes the expected eigenvalue distribution of random matrices with
i.i.d.\ entries:
\[
  \rho_{\text{MP}}(\lambda) = \frac{1}{2\pi \sigma^2 \gamma}
  \frac{\sqrt{(\lambda_+ - \lambda)(\lambda - \lambda_-)}}{\lambda}
\]
where $\lambda_\pm = \sigma^2(1 \pm \sqrt{\gamma})^2$ and $\gamma$ is the
aspect ratio of the matrix. Fourth, eigenvalues exceeding the Tracy-Widom
threshold---the upper edge of the Marchenko-Pastur bulk plus a correction
term derived from Tracy-Widom statistics---are flagged as anomalous.

This approach builds on the WeightWatcher framework~\cite{martin2021},
which demonstrated that spectral properties of weight matrices correlate
with model quality. We adapt this insight to security: a backdoor implanted
via fine-tuning or weight perturbation introduces structured spectral
signatures that deviate from the Marchenko-Pastur baseline. The spectral
watchdog detects these deviations without requiring knowledge of the
specific backdoor mechanism.

\paragraph{M4: Activation Divergence.}
Activation divergence monitoring operates on the model's internal
representations rather than its outputs. For 1\% of incoming requests
(selected uniformly at random), \mire captures an activation fingerprint:
a compressed representation of the model's intermediate activations at
designated probe layers. The fingerprint is compared against a deployment
baseline---established during model validation---using cosine similarity:
\[
  d(a, a_{\text{base}}) = 1 - \frac{a \cdot a_{\text{base}}}
  {\|a\| \, \|a_{\text{base}}\|}
\]
A divergence exceeding the empirically calibrated threshold $\tau$ triggers
an alert. This mechanism catches three classes of attack that may evade
output-level monitoring: weight replacement (wholesale substitution of model
parameters), fine-tuning attacks (gradual parameter modification via
continued training on adversarial data), and model drift (slow degradation
over time due to environmental distribution shift). The overhead is
approximately 5\,ms on sampled requests and 0\,ms on unsampled requests,
yielding an amortized cost of ${\sim}0.05$\,ms per request.

\paragraph{M5: Negative Selection Detectors.}
The fifth mechanism draws on Forrest et al.'s~\cite{forrest1994} negative
selection algorithm, originally inspired by the biological immune system's
T-cell maturation process. In the thymus, immature T-cells that react to
self-antigens are deleted; the surviving repertoire reacts only to
non-self---i.e., to foreign pathogens. \mire adapts this principle to
model behavior monitoring.

During a calibration phase, a large population of random detectors is
generated, where each detector is a predicate over model behavioral
features (output distributions, token probabilities, response latencies,
tool call patterns). Detectors that match any observed clean behavior
during calibration are deleted. The surviving detector set, by construction,
matches \emph{only} behavioral patterns not present in the clean baseline.
This is a fundamentally different approach from signature-based detection:
negative selection does not require knowledge of what a trojan looks like---it
requires only knowledge of what \emph{clean} looks like. Any behavioral
pattern sufficiently dissimilar from the clean baseline triggers a surviving
detector, regardless of the attack mechanism.

\paragraph{M6: Capability Sandbox.}
The sixth mechanism provides architectural containment---the last line of
defense when M1--M5 fail to prevent or detect backdoor activation. The
capability sandbox ensures that even a fully activated backdoor cannot
translate into meaningful impact, because the model's architectural reach
is constrained by construction:

\begin{itemize}[leftmargin=*]
\item The model cannot access external tools directly; all tool invocations
      are mediated by the L2 capability proxy, which enforces allowlists
      and rate limits independently of the model's intent.
\item The model cannot observe raw user data; all input passes through
      \pasr, which strips provenance-violating content before the model
      receives it.
\item The model cannot issue actions without authorization from the \tsa
      layer, which enforces temporal safety constraints on capability
      trajectories.
\item The model cannot exceed the capability labels assigned by \cafl,
      which bounds the maximum privilege level attainable in any session.
\item The model's behavioral trajectory is continuously monitored by \gps,
      which detects deviations from the expected capability envelope in
      real time.
\end{itemize}

\noindent
The capability sandbox embodies the core \mire philosophy: the backdoor may
\emph{activate}, but it has \emph{achieved nothing}. The model's output is
validated (M1), its behavior is monitored (M2), its weights are inspected
(M3), its activations are fingerprinted (M4), its anomalies are detected
(M5), and its architectural reach is bounded (M6). Containment does not
require detection; it requires only that every path from compromise to
impact passes through at least one mechanism that the attacker does not
control.

\subsection{Defense-in-Depth Cascade}

The six mechanisms are not independent alternatives; they form a
defense-in-depth cascade in which each layer catches threats that
penetrate the preceding layers. Consider the sequence of events when a
backdoor activates during inference:

\begin{enumerate}[leftmargin=*]
\item \textbf{Backdoor activates.} The model produces adversarial output
      in response to a trigger input.
\item \textbf{M1 blocks output.} The Output Envelope Validator rejects the
      response if it violates schema, capability, policy, or content
      constraints. For many backdoor payloads (e.g., exfiltration attempts,
      unauthorized tool calls), this is the terminal event.
\item \textbf{M2 detects behavioral drift.} Even if the adversarial output
      superficially passes the envelope validator, canary probes detect
      the statistical shift in model behavior caused by the active backdoor.
\item \textbf{M4 flags activation divergence.} On sampled requests, the
      activation fingerprint diverges from the deployment baseline,
      providing an independent signal that the model's internal state has
      been perturbed.
\item \textbf{M5 detects the anomaly.} Negative selection detectors fire on
      the novel behavioral pattern introduced by the backdoor, which by
      definition was not present in the clean calibration set.
\item \textbf{M6 contains architecturally.} Regardless of whether M1--M5
      detect the compromise, the capability sandbox ensures that the model
      cannot reach external tools, raw data, or privileged actions without
      passing through enforcement points outside its control.
\end{enumerate}

\noindent
The cascade ensures that an attacker must simultaneously evade all six
mechanisms to achieve impact---a composition of independent failure
probabilities that rapidly approaches zero.

\subsection{MIRE vs.\ BFT Consensus}

An alternative approach to model compromise is Byzantine Fault Tolerant
(BFT) consensus: run multiple independent models in parallel, compare
their outputs, and flag disagreements. While BFT consensus provides strong
guarantees in distributed systems, it is poorly suited to the LLM backdoor
problem. Table~\ref{tab:mire-vs-bft} compares the two approaches.

\begin{table}[t]
\centering
\small
\caption{Comparison of \mire containment with BFT consensus for LLM
         model compromise defense.}
\label{tab:mire-vs-bft}
\begin{tabular}{@{}lp{2.5cm}p{2.8cm}@{}}
\toprule
\textbf{Aspect} & \textbf{BFT Consensus} & \textbf{\mire} \\
\midrule
Compute cost        & $3$--$5\times$ ($\geq 3$ models)
                    & ${\sim}1.02\times$ \\
Latency overhead    & $+500$\,ms--$2$\,s
                    & $+0$--$5$\,ms \\
Compromise handling & ${\sim}50\%$ detection
                    & ${\sim}95\%$ containment \\
Same-arch.\ backdoor & Fails (correlated)
                    & Works (model-agnostic) \\
Paradigm            & Detect via disagreement
                    & Contain via architecture \\
\bottomrule
\end{tabular}
\end{table}

\noindent
The fundamental weakness of BFT consensus in the LLM setting is
\emph{correlated failure}: if all models share the same architecture,
training pipeline, or fine-tuning data, a supply-chain backdoor affects
all replicas simultaneously, and consensus yields unanimous agreement on
the adversarial output. \mire avoids this failure mode entirely because
its mechanisms are model-agnostic---they operate on outputs, activations,
weights, and architectural constraints rather than on inter-model
agreement. The ${\sim}1.02\times$ overhead of \mire (dominated by the 1\%
activation sampling in M4) is negligible compared to the $3$--$5\times$
cost of maintaining multiple model replicas for BFT voting.
